\section{Introduction}

A puzzle hunt presents its participants with a peculiar evaluative duality: the answer to
any given puzzle is either correct or it is not, but the puzzle itself is either elegant or
it is not, and that second judgment---the one that separates a puzzle worth solving from one
that merely has a solution---is far harder to make and far harder to teach.

\subsection{The Problem: Domain Expertise Is Scarce}

Evaluating creative work well requires more than intelligence and goodwill. It requires the
accumulated vocabulary of a practice: the ability to name what is wrong, to point to the
specific mechanism that fails, to distinguish a problem of clarity from a problem of
elegance from a problem of fairness. In creative domains with recognized expert communities,
this vocabulary is hard-won---it takes years of engagement with a tradition to develop the
evaluative lens that allows a practitioner to say not just ``this puzzle is unsatisfying''
but ``the confirmation mechanic is ambiguous, which shifts the cognitive load from
\emph{aha} to \emph{relief}, and that is a fairness failure, not a difficulty setting.''
Such expertise is scarce, expensive to access, inconsistently applied, and largely
unavailable at the scale required for iterative creative pipelines. The consequence is that
early-stage creative work---when the most is still malleable---is typically evaluated by
creators themselves, by generalist colleagues, or not at all. Domain expert review arrives
late, if at all, and is difficult to incorporate systematically.

\subsection{The Insight: Persona Depth Is the Variable That Matters}

Large language models have demonstrated capacity to simulate domain knowledge across a wide
range of fields~\cite{park2023generative}, and recent work on AI-simulated review has
shown that LLMs can approximate expert academic judgment under controlled
conditions~\cite{shah2023llm,dellaLibera2025profiles}. However, the standard approach to
persona simulation---supplying a brief role description such as ``you are an experienced
puzzle designer''---produces correspondingly generic output. A bare prompt elicits the
evaluative vocabulary the model associates with the role label, not the specific,
idiosyncratic frameworks that distinguish one expert's perspective from another's. The
result is feedback that is superficially domain-appropriate but operationally thin: it
identifies that problems exist without naming the mechanisms that produce them, and it
generates recommendations too vague to act on.

The insight driving this work is that the quality of AI-simulated expert feedback is not a
fixed property of the underlying model---it is a function of the richness of the persona
description. When a reviewer profile captures not just a role label and a list of
interests but a full evaluative lens, including characteristic frameworks, named concerns,
design philosophy, and the vocabulary the reviewer reaches for when a puzzle misfires, the
resulting simulated feedback is qualitatively different: more specific, more transferable,
and more actionable. The profile is not window-dressing---it is the mechanism by which the
model instantiates a distinctive evaluative stance rather than averaging across all
evaluative stances it has encountered in training.

\subsection{Our Approach: Persistent Profiles and Structured Evaluation}

We operationalize this insight through a system of \emph{persistent expert profiles}:
structured markdown documents of 80--120 lines each that capture a simulated reviewer's
design philosophy, evaluative lens, characteristic concerns, analytical vocabulary, and
voice. Profiles are loaded once per review session and reused across puzzles, building on
the token-efficient profile architecture developed in prior work~\cite{dellaLibera2025profiles}.
A panel of three reviewers is assembled from the profile library for each puzzle under
evaluation, and each reviewer scores the puzzle on a six-dimension rubric (Clarity,
Solvability, Elegance, Reading Reward, Fun, Confirmation) with a 1--5 scale per dimension,
yielding a total out of 30 with a passing threshold of 22. Scores are accompanied by
written feedback and a set of specific, actionable fix recommendations.

We evaluate this system empirically against two lower-profile conditions: a bare-prompt
baseline with no reviewer persona, and a generic profile set drawn from an earlier 12-profile
library~\cite{dellaLibera2025profiles}. The comparison isolates the effect of profile depth
from the effect of having a profile at all, allowing us to attribute outcome differences
specifically to domain specificity rather than to the presence of any persona description.
The testbed is the Age of Empires puzzle hunt, a set of original puzzles designed around
the thematic and mechanical universe of the Age of Empires game franchise. All three
puzzles evaluated in the main experiment were developed independently of the reviewer
profiles, ensuring that the profiles do not encode any information about the specific
artifacts being evaluated.

\subsection{What Is a Puzzle Hunt?}

For readers unfamiliar with the domain: a \emph{puzzle hunt} is a competitive or
collaborative event in which teams solve a sequence of puzzles, each yielding an
answer---typically a word or phrase---that feeds into a larger \emph{meta-puzzle} resolving
the hunt's central mystery. The puzzles themselves range across a wide variety of types
(word puzzles, logic grids, cryptic crosswords, visual puzzles, knowledge challenges) and
are evaluated against a demanding set of craft criteria: they must be unambiguously solvable
by the target audience, must reward careful reading, must produce an \emph{a-ha} moment at
the point of solution, and must feel elegant and fair in retrospect. The hunt community has
developed a sophisticated evaluative vocabulary around these criteria, and recognized
practitioners---editors at MIT Mystery Hunt, MITMH-adjacent events, and independently
organized hunts---bring a distinctive and largely codified expert perspective to the
work~\cite{colton2012computational}. This makes puzzle hunt design an unusually appropriate
testbed for studying AI-simulated expert review: the expertise is real, specialized, and
reasonably well-documented, but access to it is constrained by the small size and
geographic concentration of the expert community.

\subsection{Contributions}

This paper makes the following contributions:

\begin{enumerate}

\item \textbf{The calibration inversion.} Rubric-only evaluation (no reviewer identity)
inverts the quality tier hierarchy: across 15 real published puzzles, it produces an 80\%
pass rate for Standard Hunt puzzles but only 40\% for MIT/Elite-tier puzzles. The mechanism
is traceable: elite puzzles deliberately sacrifice Clarity and Solvability for Elegance and
Reading Reward, and an equal-weight rubric treats that trade-off as a defect. This finding
is a concrete warning for practitioners---rubric-only evaluation appears objective but is
miscalibrated for ambitious creative work.

\item \textbf{Profile-based AI expert panel system.} A complete system for assembling and
deploying AI reviewer panels for creative work evaluation, built on persistent structured
profiles encoding domain-specific evaluative frameworks. The system is described in
sufficient detail to support replication and adaptation to other creative domains.
\emph{Note: all quantitative findings are valid within the authors' constructed evaluation
framework; the rubric, tier assignments, and thresholds were constructed by the research
team and have not been externally calibrated against a held-out human expert standard.}

\item \textbf{Empirical eight-condition profile-depth ablation.} A controlled experiment
across 15 real published puzzles from MIT Mystery Hunt 2023, Huntinality, and Teammate Hunt,
comparing conditions from bare prompting to full domain-specific profiles to profile plus
principles. The ablation isolates the contribution of domain specificity, profile
component type, and rubric design.

\item \textbf{Taxonomy of surfaced analytical frameworks.} Six named evaluative frameworks
(\emph{world-model consistency}, \emph{a-ha economy}, \emph{load-bearing test},
\emph{social fabric}, \emph{perceptual-shift}, \emph{visual grammar}) that appear in
C6 reviews but in no input text provided to the reviewers, constituting a contribution to
the vocabulary of puzzle design evaluation.

\item \textbf{Open profile library and replication package.} The full set of 29
domain-specific reviewer profiles (C8 format, $\sim$75 lines), the seven-dimension weighted
rubric, all review transcripts, and analysis scripts, released to support replication and
adaptation to other creative evaluation domains.

\end{enumerate}

\subsection{Paper Organization}

Section~\ref{sec:related} situates this work relative to AI-simulated review, persona-based
prompting, computational creativity evaluation, and the emerging literature on puzzle design
as a research domain. Section~\ref{sec:methodology} describes the profile architecture,
scoring rubric, and experimental design in detail. Section~\ref{sec:evaluation} reports
results from the profile ablation and cross-hunt generalizability experiments.
Section~\ref{sec:discussion} interprets the results and discusses limitations.
Section~\ref{sec:conclusion} summarizes contributions and directions for future work.
